%%=====================================================================
%%  1 前言
%%  【规范】一级标题：小 4 号黑体，顶格；正文：小 4 号仿宋，首行缩进 2 字符。
%%=====================================================================
\section{前言}

随着互联网和移动设备信息技术的快速发展，智能移动终端设备得到了极大的普及，人们获取外界信息的习惯也从曾经的报纸、书籍慢慢转变为通过互联网直接获取相关信息。众所周知，学习习惯对于人们来说是需要慢慢培养的，因此伴随着互联网的高速发展，人们从传统的纸质书籍学习拓展到了如今的线上学习，这是一个伟大的转变。相较于传统的学习方式，线上学习具有以下优势。

首先，不受学习时间和空间的限制。许多上班族都会选择网上学习的方式学习一些跟自己职业生涯相关的知识，特别是对于部分职场新人，由于个人工作能力的欠缺、工作经验的不足，往往会通过线上学习的方式获取需要的专业知识。对于大部分上班族来说，白天需要专心于公司的工作，而晚上又可能需要处理各种各样的个人问题或者家庭问题，导致没有办法挤出大量的时间去实体班系统化地学习。但是，线上学习则可以方便地配合个人的时间与地点，只要有网络和移动终端设备就可以随时随地利用个人碎片化的时间快速开始学习。

其次，线上学习费用更低。自 2005 年以来，各种类型的线上教育平台如雨后春笋般涌现，其发展如此迅猛的原因之一就是费用和资源得到了大量节省。相对于传统教育的模式，线上教育节省了住宿、教师、教室、资料等大量资源，据不完全统计，线上学习大致花费仅相当于传统现场学习费用的 30\%---50\%\footnote{该比例引自在线教育行业公开调研数据，不同机构的统计口径存在差异。}，却可以达到相差无几的学习效果\cite{bib:moodle}。

最后，更强的可重复性与个性化学习。线上学习可以根据个人学习需求和学习进度反复学习，重复未理解的部分，以便更好地掌握所学内容并充分巩固学习效果，避免在传统课堂学习中容易出现的"学过就忘"的问题。同时，线上学习能够较好地实现个性化学习，人们可以根据个人的时间安排，依照个人的需求、文化背景、学习习惯来选择学习内容，有效地增强学习的针对性和提高学习效率。

在移动终端上，一个 APP 应用程序需要 Android 和 iOS 两套技术代码，应用程序的上传步骤繁琐，开发周期较长。同时，一个优秀的 APP 应用的研发与推广，不仅需要配备优秀的产品经理、市场运营人员，还需要高额的推广成本。此外，部分 APP 界面操作复杂，占用较多手机内存，需要用户授予诸多权限等因素，造成许多用户不愿意轻易下载 APP 应用程序。

微信小程序不需要用户安装，即开即用，节省了许多用户流量和安装时间。相较于原生 APP，小程序不但在用户使用体验感上差别不大，而且在开发成本、获客成本、下载便捷程度以及轻应用方面具有更大的优势。微信小程序与第三方应用不同，无需频繁地跳出微信再重新回到微信，减少了用户许多重复性的操作，提升了用户体验。微信小程序依附于微信生态系统的环境下，因而后期的技术维护比原生 APP 更加容易。截止到 2016 年第二季度，微信已经覆盖国内 94\% 以上的智能移动终端，月活跃用户量达到 8 亿以上，用户覆盖超过 200 个国家、20 种语言，微信公众账号的总数超过 800 万个，移动应用程序的对接数量超过 85000 个\cite{bib:xiong}。

这次我们所设计的课题系统，正是移动终端设备社交平台与社会生活之间紧密关系的体现。随着微信的迅猛发展，微信逐渐成为现代人生活中一个重要的组成部分，它对用户信息的获取、电子商务、生活服务、即时交流等许多方面产生了深刻的影响。微信小程序依托微信庞大的用户流量，以一种全新的方式连接用户与服务。本论文以增强移动端用户 PMP 考证刷题的便捷性和实时性为目的，设计并实现了基于微信小程序的 PMP 刷题系统。
